\section{Overland Movement}\label{Overland Movement}
  This section provides rules governing overland movement.
  Not every game should think about overland movement travel speed in a detailed way.
  It's fine to just spend ``a few days'' walking around between various important locations.
  However, sometimes the details are important, such as when you are on a strict timetable.
  The GM can tell you when overland movement matters.

  \subsection{Standard Travel Days}
    Characters covering long distances cross-country use overland movement.
    Overland movement is measured in miles per hour or miles per day.
    A day normally represents 8 hours of actual travel time.
    However, sailing ships and other methods of travel that keep moving without requiring a rest are listed with a full 24 hours of travel time.

    You can make an Endurance check to push beyond a standard 8-hour travel day.
    In addition, you can make an Endurance check to travel faster within a normal travel day.
    For details, see \pcref{Endurance}.

    Standard travel distances on foot are listed in \tref{Travel Distance By Movement Speed}.
    When using mounts or ships, \tref{Mounts and Vehicles} will be more convenient.

    \begin{columntable}
      \columncaption{Travel Distance By Movement Speed}
      \begin{dtabularx}{\columnwidth}{>{\lcol}X c c c c}
        & \multicolumn{4}{c}{\tdash\tdash\tdash Speed \tdash\tdash\tdash} \tableheaderrule
        & 15 feet    & 20 feet  & 30 feet  & 40 feet  \\
        One Hour (Overland) &            &          &          &          \\
        Walk                & 1-1/2 mile & 2 miles  & 3 miles  & 4 miles  \\
        Hustle (Exertion)   & 3 miles    & 4 miles  & 6 miles  & 8 miles  \\
        One Day (Overland)  &            &          &          &          \\
        Walk                & 12 miles   & 16 miles & 24 miles & 30 miles \\
        Hustle (Exertion)   & \tdash     & \tdash   & \tdash   & \tdash   \\
      \end{dtabularx}
    \end{columntable}

    \begin{columntable}
      \columncaption{Mounts and Vehicles}
      \begin{dtabularx}{\columnwidth}{>{\lcol}X l l}
        \tb{Mount/Vehicle}                         & \tb{Per Hour} & \tb{Per Day} \tableheaderrule
        Mount (carrying load)                      &             &          \\
        \tind Light horse or warhorse              & 5 miles     & 40 miles \\
        \tind Draft horse                          & 4 miles     & 32 miles \\
        \tind Pony or war pony                      & 4 miles     & 32 miles \\
        \tind Donkey or mule                       & 3 miles     & 24 miles \\
        \tind Dog, riding                          & 4 miles     & 32 miles \\
        \tind Cart or wagon                        & 2 miles     & 16 miles \\
        \tb{Ship}                                  &             &          \\
        \tind Raft or barge (poled or towed)\fn{1} & 1/2 mile    & 4 miles  \\
        \tind Keelboat (rowed)\fn{1}               & 1 mile      & 8 miles \\
        \tind Rowboat (rowed)\fn{1}                & 1-1/2 miles & 12 miles \\
        \tind Sailing ship (sailed)                & 2 miles     & 48 miles \\
        \tind Warship (sailed and rowed)           & 2-1/2 miles & 60 miles \\
        \tind Longship (sailed and rowed)          & 5 miles     & 120 miles \\
        \tind Galley (sailed and rowed)            & 6 miles     & 144 miles \\
      \end{dtabularx}
      1 Rafts, barges, keelboats, and rowboats are used on lakes and rivers.
      If going downstream, add the speed of the current (typically 3 miles per hour) to the speed of the vehicle. In addition to 10 hours of being rowed, the vehicle can also float an additional 14 hours, if someone can guide it, so add an additional 42 miles to the daily distance traveled. These vehicles can't be rowed against any significant current, but they can be pulled upstream by draft animals on the shores.
    \end{columntable}

  \subsection{Overland Terrain}
    Travelling over a flat, paved highway is much faster than trailblazing through a jungle.
    You can use \tref{Terrain and Overland Movement} as a reference for common terrain.

    A highway is a straight, major, paved road.
    A road is typically a dirt track.
    A trail is like a road, except that it allows only single-file travel and does not benefit a party traveling with vehicles.
    Trackless terrain is a wild area with no significant paths.

    \begin{columntable}
      \columncaption{Terrain and Overland Movement}
      \begin{dtabularx}{\columnwidth}{>{\lcol}X c c c}
        \tb{Terrain}   & \tb{Highway} & \tb{Road or Trail} & \tb{Trackless} \tableheaderrule
        Desert, sandy  & \mult1       & \mult1/2           & \mult1/2 \\
        Forest         & \mult1       & \mult1             & \mult1/2 \\
        Hills          & \mult1       & \mult3/4           & \mult1/2 \\
        Jungle         & \mult1       & \mult3/4           & \mult1/4 \\
        Moor           & \mult1       & \mult1             & \mult3/4 \\
        Mountains      & \mult3/4     & \mult3/4           & \mult1/2 \\
        Plains         & \mult1       & \mult1             & \mult3/4 \\
        Swamp          & \mult1       & \mult3/4           & \mult1/2 \\
        Tundra, frozen & \mult1       & \mult3/4           & \mult3/4
      \end{dtabularx}
    \end{columntable}

